\documentclass[11pt, letterpaper,titlepage]{article}

% --- Essential Packages ---
\usepackage[utf8]{inputenc}
\usepackage[margin=1in]{geometry}
\usepackage{fancyhdr} % For custom headers and footers
\usepackage{graphicx}
\usepackage{amsfonts}

\usepackage{amsmath}
\usepackage{float}

\usepackage{hyperref}
\usepackage{setspace}

\usepackage{ragged2e}

% --- Header/Footer Configuration ---
\pagestyle{fancy}
\fancyhf{} % Clear all default header and footer settings
\fancyhead[L]{DIGITALLY CONTROLLED ANALOG MIXER} % Left header text
\fancyhead[C]{} % Center header text (empty)
\fancyhead[R]{\thepage} % Right header (page number)
\renewcommand{\headrulewidth}{0.4pt} % Line below header
\linespread{1.2}





\begin{document}
\begin{titlepage}
    \centering
    \vspace*{1cm}
    
    {\huge \textbf{Analog Mixing Console With Digital Control\\ \vspace{0.5cm}and Automated Signal Path Configuration}} \\[1cm]
    
    \vspace{0.5cm}
    
    {\huge Capstone Project Proposal}
    
    \vspace{5cm}
    




    
    {\LARGE \textbf{Skylar Denno}}
    
    \vspace{0.3cm}
    
    {\LARGE \textbf{Evan Forcucci}}
    
        \vspace{0.3cm}
    
    {\LARGE \textbf{Daniil Petrov}}
    
        \vspace{0.3cm}
    
    {\LARGE \textbf{Xinhao Chen}}
    
        \vspace{0.3cm}
    
    {\LARGE \textbf{Daniel Salas}}
    
        \vspace{0.3cm}
    
    {\LARGE \textbf{Luke Garrity}}
    

    
    \vspace{1cm}
    
        {\Large \textbf{Advisor: }Milad Siami }
    
        \vspace{2cm}

    
        {\large \textbf{Northeastern University}} \\
        
    Department of Electrical and Computer Engineering \\
    \vspace{1cm}
    
    
    \today

\end{titlepage}

\justify

\tableofcontents

\section{Abstract}



\section{Introduction}


\section{System Explanation}

\subsection{Fundamental Design and Vision}

\subsection{Analog System Overview}

\subsection{STM32 Microcontroller and Digital Overview}


\section{Technical Details}




\subsection{Voltage Controlled Amplification}
\label{sssec:vca}

To remotely control parameters of the analog signal processing, all the required processing may be (and usually is, even in fully analog systems) designed to only require variable resistors and variable gain. For example, to change the resonant frequency of a band on the equalizer, we may keep the capacitor value constant and vary only the resistor to change the resonant frequency. There are several ways to implement voltage controlled resistors. We plan to use operational transconductance amplifiers (OTAs) to implement this voltage controlled resistance, specifically the LM13700, for its output Darlington buffers and linearizing diodes, which help it achieve a THD of $<1\%.$ To implement the variable gain, voltage controlled amplifiers (VCAs) will be used.

The first place that gain must be remotely controlled is in the input section, which takes in either a balanced line level (preamplified) signal, or a balanced quiet microphone signal. For line inputs, an INA2137 line input receiver with fixed unity gain will impedance balance and convert the double ended signal into a single ended signal for processing. A VCA with a control voltage will be used to vary the gain. For a microphone input, a PGA2500 digitally controlled microphone preamplifier will be used on an SPI bus to allow control of the microphone gain.


The second place gain will need control is on the auxiliary send section, where variable amounts of audio from each channel may be summed to auxiliary signal buses for group external processing. These busses will be fed by VCAs, specifically the SSI2164. The third place for voltage controlled gain is the feed into the main summing bus. For this, two VCAs (SSI2164s) will be used per channel. One VCA will feed the left channel summing bus, and the other VCA will feed the right channel summing bus. A pan (stereo panorama balance) knob will control the difference in gain between the two VCAs for each channel, and the channel fader (volume control) will scale both VCAs equally. 

\subsubsection*{Component: PGA2500}

The PGA2500 digitally controlled microphone preamplifier accepts its gain setting as a number, so it needs no control DAC. A
16-bit control word travels over a serial port using the SPI protocol. The gain range runs from
10\,dB to 65\,dB in 1\,dB steps, with a separate unity gain setting. Several
devices may share one port in a daisy chain, which keeps the wiring count low across a
large number of channels. The device also provides four programmable digital outputs. These
outputs drive the pad, the phantom power supply (for condenser / capacitance microphones), the high-pass filter, and the
polarity reverse relay.

The zero-crossing detector delays each gain change until the input signal
crosses zero. This limits the glitch energy from the switched gain network, preventing pops and clicks on the output. A
channel with no input signal therefore shows some delay before the new gain
takes effect. Preamplifier recall is exact. A switched resistor network sets
the gain, and no control voltage takes part.


\subsubsection*{Component: SSI2164}
The SSI2164 provides four VCAs in one package. Each VCA takes a current input,
produces a current output, and accepts a ground-referenced control port. The
control law is exponential:
%
\begin{equation}
    V_C = -\left(33\,\mathrm{mV/dB}\right) G ,
    \label{eq:vca-law}
\end{equation}
%
where $G$ is the gain in decibels. A control span of $-660$\,mV to $+3.3$\,V
therefore covers the full range of $+20$\,dB to $-100$\,dB.

Equation~\eqref{eq:vca-law} gives a useful property for recall. One control DAC
step always moves the gain by the same number of decibels. A 12-bit control DAC
across a 4.096\,V span gives a 1\,mV step, or 0.03\,dB. This resolution holds
at the top of the fader and at 90\,dB of attenuation.

Noise on the control port becomes gain modulation, so the control path needs
care. Noise of 1\,mV RMS produces sidebands about 55\,dB below the signal. A
target of 10\,$\mu$V RMS in the audio band moves those sidebands below
$-95$\,dBc. The control DAC, its voltage reference, and its output buffer must
all meet this target.

Each control DAC step also makes a step in gain. A first-order RC filter on
each control line removes the glitch energy from the step. The corner frequency
cannot go very low, because the fader must still respond fast. We set the
corner near 100\,Hz and update the control DAC at 1\,kHz. The control processor
interpolates fader moves, so the hardware receives many small steps instead of
one large step. The control voltage settles within about 15\,ms after the last
update.

\subsection{Voltage Controlled Resistors}

Parts of the channel need a variable resistance and not a variable gain. The
equalizer corner frequencies and the filter $Q$ are the main examples. An
LM13700 OTA supplies this resistance. Transconductance follows the amplifier
bias current:
%
\begin{equation}
    g_m = \frac{I_\mathrm{ABC}}{2 V_T} ,
    \qquad
    R_\mathrm{eq} = \frac{1}{g_m} = \frac{2 V_T}{I_\mathrm{ABC}} ,
    \label{eq:ota-gm}
\end{equation}
%
where $V_T = kT/q$ is the thermal voltage. A bias current from 1\,$\mu$A to
500\,$\mu$A gives a resistance from about 51\,k$\Omega$ down to about
100\,$\Omega$. Three decades of range cover a wide sweep of one corner
frequency.

The OTA stays linear only for small differential input voltages. A resistive
divider holds the signal at the OTA input below about 20\,mV peak. The
linearizing diodes extend this limit. Both methods cost signal-to-noise ratio,
so the divider ratio sets the noise floor of the equalizer stage.

An op-amp voltage-to-current converter derives $I_\mathrm{ABC}$ from the
control DAC. This keeps the bias current independent of the supply rail.
Transconductance follows bias current in a linear way, so corner frequency also
follows the control DAC code in a linear way. A lookup table in the control
processor maps a target frequency to the correct code. The analog circuit
therefore needs no exponential converter.

\subsubsection*{Component: LM13700}

The LM13700 is a dual operational transconductance amplifier, manufactured originally by National Semiconductor, and now TI. Unlike an op-amp (which is a voltage controlled voltage source), OTAs produce an output CURRENT proportional to the differential input voltage, just like a resistor, and the scaling factor (the effective resistor value) is set by an external amplifier bias current: $gm = I_ABC / 2V_T$, or roughly 19 mS per mA of bias at room temperature. That makes it a general-purpose analog multiplier building block, which is why it turns up as a VCA, voltage-controlled filter, voltage-controlled resistor, or oscillator, and why it became a staple of synthesizer design.

Each of the two channels adds linearizing diodes and a Darlington buffer. The diodes matter because the bare differential pair is linear only over a few tens of millivolts, so most designs either attenuate hard at the input or use the diodes to trade input range against noise. The buffers give you a high-impedance follower to drive the next stage without loading the current output. It runs on anything from about ±2 V to ±18 V, with bias currents up to a couple of milliamps.

\subsection{Calibration and Temperature Sensitivity}
\label{sssec:temperature}

The gain constant in Equation~\eqref{eq:vca-law} and the thermal voltage in
Equation~\eqref{eq:ota-gm} both scale with absolute temperature. This is the
largest threat to recall accuracy in the channel.

The gain constant of the SSI2164 rises with absolute temperature. A rise of
30\,$^\circ$C moves the constant from about 33\,mV/dB to about 36\,mV/dB. A
control voltage stored for 40\,dB of attenuation then produces about 36\,dB.
The error grows with distance from unity gain and falls to zero at 0\,dB. The
thermal voltage $V_T$ drifts in the same manner. The equivalent resistance
rises by about 0.33\,\% per $^\circ$C, and any corner frequency set by that
resistance falls by the same fraction.

We correct this drift in two places. A resistor with a $+3300$\,ppm/$^\circ$C
coefficient scales the control voltage at each SSI2164 control port. The
control processor also reads a temperature sensor on each module and trims the
control DAC codes. The sensor correction handles the residual error and the
slow drift of the voltage reference.

Device tolerance also limits recall accuracy. Each module receives a
calibration table at the end of manufacture. The test system measures gain
against control DAC code at several points. It stores the corrections in
nonvolatile memory on the module. The control processor reads this table at
power-up and applies it to every recall. The table travels with the module, so
a board swap does not change the calibration. This recall accuracy however, is still likely to be more precise than the manual recall of physically turning potentiometers to the positions indicated on a recall sheet.



\subsubsection*{Recall Sequence}
\label{sssec:recall}

Reconfiguring the console's state follows a fixed order. This order prevents transients at the output.

\begin{enumerate}
    \item Ramp the channel VCA to full attenuation.
    \item Write the routing, pad, polarity, and phantom power states.
    \item Write the gain word to the PGA2500.
    \item Write the control DAC codes for the voltage-controlled resistors.
    \item Wait for the control filters to settle.
    \item Ramp the channel VCA to the stored control voltage.
\end{enumerate}


\subsection{State-Variable Filter Equalizer}

\subsection{Dynamic Range Compressor}

\subsection{Summing Bus}

\subsection{Physical Controls and Knob Position Indication}

\subsection{Motorized Faders}

\subsection{LCD Displays}


\label{ssec:displays}

The knob rings, the motorized fader, and the button lamps already report
position and state. They cannot report a name, and they cannot report a shape.
A small display adds these two things. Each channel strip carries one Adafruit
4421 module, which holds a 1.54\,inch IPS panel of 240 by 240 pixels and an
ST7789 controller. The screen shows the track name, the transfer curve of the
equalizer, and the gain reduction of the compressor. Adafruit supplies this
part as a bare module on a flexible cable, so each strip also needs a small
connector for the 22 pin "flat-flex" ribbon cable.

One controller board drives four screens on a single SPI bus configured as transmit only (no data return from screens to controller). The four modules will 
share the clock line and the data line (MISO). Each module receives its own chip select line. Two pins on the controller will be used to select the module to talk to using a two bit number. A 2-to-4 decoder generates the four chip select lines and saves two pins on the processor. The backlight of all four displays will be controlled by another pin, and pulse width modulation (PWM) will be used to vary its brightness.

Bus bandwidth sets the whole update strategy. A full frame holds 57,600 pixels
at 16 bits per pixel, or 115,200 bytes. A clock of 40\,MHz therefore needs
23\,ms for one full frame, and 92\,ms for all four screens. A screen that always
redraws in full is too slow for a meter.

The ST7789 holds its own frame memory. The controller sets an address window
and then writes only the pixels inside that window. The design uses this
feature and divides each screen into three zones.

\begin{itemize}
    \item The track name changes only on a rename or a snapshot recall. The
          controller redraws this zone on that event.
    \item The equalizer curve changes only when the user moves an encoder. One
          hand moves one encoder, so at most one channel redraws this zone at
          a time.
    \item The gain reduction meter changes all the time, so this zone stays
          small. A bar of 16 by 200 pixels needs 6400 bytes, or 1.3\,ms. Four
          meters at 30 updates per second use about 15\,\% of the bus.
\end{itemize}

A full frame buffer for four screens would need 461\,kB of memory, and the
controller does not hold one. The controller renders one tile at a time into a
small buffer instead. A tile of 240 by 32 pixels needs 15\,kB. The controller
keeps two tiles and uses DMA. The processor draws into the second tile while
the DMA engine sends the first tile to the bus. It is likely that during design, we will sacrifice colors to reduce the bits per pixel, in order to increase framerate. We also plan to have a single, larger display in the center / master section, for the user to configure settings and save and load console states.

\subsection{GPIO Expansion}

\subsection{Control Voltage DACs}
\label{ssec:cv-dacs}

Each console channel needs many control voltages at once. These are the VCA control
ports and the equalizer corner frequencies. An STM32G474 supplies these levels
through four independent I\textsuperscript{2}C buses. Each bus carries one
console channel, so one processor serves four channels.

Each bus holds eight devices. Three MCP23017 expanders provide the switch and
lamp lines described above. Five MCP4728 converters provide the control
voltages. The MCP4728 holds four 12-bit outputs, so each bus produces 20
control voltages. One processor therefore produces 80 control voltages in
total.

Every MCP4728 leaves the factory with the same default address, so the five
devices on a bus need distinct addresses. The address bits live in on-chip
EEPROM. A programming step at the end of manufacture writes them once. This
step needs a precise transition of the LDAC pin against the I\textsuperscript{2}C
clock, so the test fixture performs it at a reduced clock rate. The LDAC pin
serves a second purpose during normal use. The processor writes all five
devices and then pulses LDAC, and all 20 outputs of a channel move together.

\subsubsection*{Bus Bandwidth}
\label{sssec:bus-bandwidth}

The bus runs in Fast mode at 400\,kHz. One bit period is therefore
2.5\,$\mu$s, and one byte takes 22.5\,$\mu$s because each byte carries an
acknowledge bit. Two transactions repeat every frame. The Fast Write command of
the MCP4728 loads all four channels in nine bytes. A two-register read of the
MCP23017 returns both ports in five bytes. Table~\ref{tab:i2c-budget} gives the
budget.

\begin{table}[t]
\centering
\caption{Bus load for one channel at a 400\,Hz frame rate. Each figure includes
an allowance for the start, stop, and bus free periods.}
\label{tab:i2c-budget}
\begin{tabular}{l r r r}
\toprule
\textbf{Transaction} & \textbf{Bytes} & \textbf{Time} & \textbf{Total} \\
\midrule
MCP4728 fast write, 5 devices & 9 & 210\,$\mu$s & 1050\,$\mu$s \\
MCP23017 port read, 3 devices & 5 & 123\,$\mu$s & 369\,$\mu$s \\
\midrule
Traffic per frame & & & 1419\,$\mu$s \\
Frame period at 400\,Hz & & & 2500\,$\mu$s \\
Bus load & & & 57\,\% \\
\bottomrule
\end{tabular}
\end{table}

The bus therefore carries the full update in 1.42\,ms of a 2.5\,ms frame. The
remaining 1.08\,ms absorbs the lamp writes, the retries, and the occasional
configuration write. A saturated bus would sustain about 700 frames per second,
so the 400\,Hz rate keeps a margin of nearly two to one. The bus stays at
400\,kHz even though the MCP4728 accepts 3.4\,MHz, because the MCP23017 shares
the same wires and stops at 1.7\,MHz. The measured rate already meets the need,
and a slower clock produces slower edges and less radiated noise.

A frame rate of 400\,Hz sets a Nyquist limit of 200\,Hz. No control voltage can
carry a component above that frequency. This limit costs nothing, because no
parameter in the channel needs to move that fast.

\subsubsection*{Reference, Span, and Output Stage}
\label{sssec:cv-output}

Each MCP4728 takes its reference from VDDA. This is a dedicated 3.3\,V linear
supply that feeds the converters and nothing else. A ferrite bead and local
decoupling isolate every device on the rail. The converter output is therefore
ratiometric to VDDA, and this suits the design. The analog-to-digital converters share the
same rail, so a measured fader position and the control voltage derived from it
scale together.

The internal 2.048\,V reference of the MCP4728 offers no advantage at this
supply voltage. It gives a full scale of only 2.048\,V at unity gain, and the
gain of two would ask for 4.096\,V from a 3.3\,V rail. A reference of VDDA
therefore gives the widest span available. A full scale of 3.3\,V across 12 bits
gives a step of 0.81\,mV at the converter output.

A TL074 stage follows each output. The stage applies a gain and an offset, and
a reference voltage at the second input sets that offset. Two resistors per
channel therefore fix the span. The widest span reaches about $\pm 13$\,V,
because the TL074 does not swing to its rails. A design that needs a true
$\pm 15$\,V output must run the stage from $\pm 18$\,V rails.

The span is not the same for every destination, and this matters. Twelve bits
spread across the full 26\,V output range give a step of 6.3\,mV, or 0.19\,dB at
a VCA control port. That step is too coarse for a fader. Destinations that need
fine resolution therefore receive a narrow span. A VCA control port spans
$-0.66$\,V to $+3.3$\,V, which needs a stage gain of 1.2. One step at that port
moves the control voltage by 0.97\,mV, or 0.03\,dB. Bipolar destinations such as a pan
control or a filter offset use the wide span, where a coarse step does no harm.

The TL074 adds an input offset voltage of a few millivolts, and the stage gain
multiplies it. On a narrow span this offset produces a fixed gain error of a
few tenths of a decibel. 
\subsubsection*{Reconstruction and Slew}
\label{sssec:cv-filter}

A capacitor across the feedback resistor of each TL074 stage forms a single
pole near 100\,Hz. This pole does three jobs. It attenuates the images that the
zero-order hold places at the frame rate. It limits the noise bandwidth of the
stage, so the amplifier contributes well under 1\,$\mu$V in the audio band. It
also slows every parameter change.

The image rejection deserves an honest figure. A single pole at 100\,Hz
attenuates a 400\,Hz component by about 12\,dB, which is modest. The step size
therefore carries the rest of the argument. The control processor limits how
far a code may move in one frame, so a large parameter change arrives as many
small steps. Large residual steps appear only during a fast deliberate move,
and the move itself masks them.

The slew limit is a choice and not only a constraint. Slower changes in parameters usually sound more natural and musical, and they prevent artifacts that may result from sudden edges in the control voltage. The filter
therefore serves the sound of the console as well as its noise budget.


\subsection{Power Supply}

\subsection{Physical Manufacturing Technicalities}

\subsection{Computer Connection}


\section{Division of Tasks}

\subsection{Digital Architecture Design}

\subsection{Firmware}

\subsection{Analog Systems Design}

\subsection{Control and UI Design}

\subsection{PCB and CAD Design}


\section{Budget Proposal}

\subsection{Component Selection}

\subsection{Component Sources}

\subsection{PCB and CAD Services}

\section{Timeline}

\subsection{Design and Prototyping}

\subsection{Implementation and Iteration}

\subsection{Final Assembly}

\section{Conclusion}




\section{References}




	
\end{document}
